% =====================================================================
%  chapters/minggu-15.tex
%  Minggu 15 — Workshop Proyek Akhir: Implementasi, Deployment &
%              Progress Clinic
% =====================================================================

\chapter{Workshop Proyek Akhir --- Implementasi, Deployment \& Progress Clinic}
\label{chap:minggu-15}

% ---------------------------------------------------------------------
% 1. CAPAIAN PEMBELAJARAN
% ---------------------------------------------------------------------
\begin{capaian}
Setelah menyelesaikan bab ini, mahasiswa diharapkan mampu:
\begin{itemize}
  \item mengimplementasikan halaman web statis (HTML + CSS + JS) sesuai \emph{mockup} Figma yang sudah dibuat;
  \item mendeplöy website ke GitHub Pages agar bisa diakses publik;
  \item menyelesaikan kendala teknis melalui sesi \emph{clinic} dengan pengajar;
  \item memastikan pembagian kerja berjalan merata dan setiap anggota memahami kode tanggung jawabnya.
\end{itemize}
\end{capaian}

% ---------------------------------------------------------------------
% 2. TUJUAN PRAKTIK
% ---------------------------------------------------------------------
\begin{tujuanpraktik}
Pada sesi praktik ini, mahasiswa akan:
\begin{itemize}
  \item menyinkronkan \emph{repository} dan memulai implementasi halaman utama;
  \item menulis struktur HTML semantik dan CSS sesuai \emph{mockup};
  \item mendeploy website ke GitHub Pages;
  \item mengikuti sesi \emph{clinic} untuk perbaikan teknis;
  \item menyiapkan bahan presentasi untuk minggu 16.
\end{itemize}
\end{tujuanpraktik}

% ---------------------------------------------------------------------
% 3. PERALATAN
% ---------------------------------------------------------------------
\section{Peralatan yang Dibutuhkan}
\label{sec:peralatan-15}

\begin{itemize}
  \item \textbf{VS Code} --- menulis HTML, CSS, JS.
  \item \textbf{Live Server} (ekstensi VS Code) --- \emph{preview} halaman secara langsung.
  \item \textbf{Browser} (Chrome/Firefox) --- \emph{testing} dan \emph{debugging}.
  \item \textbf{Figma} --- referensi desain saat coding.
  \item \textbf{Git \& GitHub} --- \emph{version control} dan deployment.
  \item \textbf{GitHub Copilot atau agentic coding tool} (opsional) --- bantuan penulisan kode.
\end{itemize}

% ---------------------------------------------------------------------
% 4. SYNCHRONISASI REPOSITORY
% ---------------------------------------------------------------------
\section{Synchronisasi Repository}
\label{sec:sync-repo}

\begin{langkah}
  \item Buka terminal di VS Code, pastikan berada di \emph{branch} yang benar:
\begin{lstlisting}[language=bash, caption={Cek branch}]
git status
git checkout final-project
\end{lstlisting}
  \item Ambil perubahan terbaru dari GitHub:
\begin{lstlisting}[language=bash, caption={Pull perubahan terbaru}]
git pull origin final-project
\end{lstlisting}
  \item Buka file \texttt{index.html} yang sudah ada (atau buat jika belum).
\end{langkah}

\begin{tangkapanlayar}
  {assets/images/minggu-15/terminal-git-status.png}
  {Terminal di VS Code --- menunjukkan git status branch final-project.}
  {Buka terminal di VS Code, jalankan \texttt{git status} dan pastikan branch sudah sesuai.}
\end{tangkapanlayar}

% ---------------------------------------------------------------------
% 5. IMPLEMENTASI HALAMAN UTAMA
% ---------------------------------------------------------------------
\section{Implementasi Halaman Utama}
\label{sec:implementasi}

Fokuskan waktu pada satu halaman utama (\emph{landing page} / \emph{homepage}). Multi-halaman bisa ditambahkan jika waktu memungkinkan.

\subsection{Struktur HTML Dasar}
\label{subsec:html-dasar}

Buka \texttt{src/index.html} dan buat struktur dasar berikut:

\begin{lstlisting}[language=HTML, caption={Struktur dasar HTML semantik}, label={lst:html-struktur}]
<!DOCTYPE html>
<html lang="id">
<head>
  <meta charset="UTF-8">
  <meta name="viewport"
        content="width=device-width, initial-scale=1.0">
  <title>[Nama Website Redesign]</title>
  <link rel="stylesheet" href="css/style.css">
</head>
<body>
  <!-- Header & Navigasi -->
  <header>
    <nav>
      <!-- Logo + menu navigasi sesuai mockup -->
    </nav>
  </header>

  <!-- Hero Section -->
  <section class="hero">
    <!-- Judul, deskripsi singkat, tombol CTA -->
  </section>

  <!-- Konten Utama -->
  <main>
    <!-- Bagian-bagian konten sesuai mockup -->
  </main>

  <!-- Footer -->
  <footer>
    <!-- Informasi kontak, copyright, tautan -->
  </footer>

  <script src="js/script.js"></script>
</body>
</html>
\end{lstlisting}

\begin{tip}
Gunakan tag HTML semantik: \texttt{<header>}, \texttt{<nav>}, \texttt{<main>}, \texttt{<section>}, \texttt{<footer>}. Hindari penggunaan \texttt{<div>} berlebihan tanpa makna. Tambahkan \texttt{lang="id"} di tag \texttt{<html>} untuk aksesibilitas.
\end{tip}

\subsection{Pembagian Kerja Kelompok}
\label{subsec:pembagian-kerja}

Agar tidak bentrok, gunakan \emph{branch} terpisah per anggota atau bagian file terpisah.

\textbf{Opsi A: Branch terpisah (disarankan):}
\begin{lstlisting}[language=bash, caption={Branch per anggota}]
# Anggota A
git checkout -b feat/header-navigasi

# Anggota B
git checkout -b feat/konten-utama

# Anggota C
git checkout -b feat/footer-responsive
\end{lstlisting}

Setiap anggota bekerja di \emph{branch} sendiri, lalu \emph{merge} ke \texttt{final-project} setelah selesai.

\textbf{Opsi B: Bagian file terpisah (lebih sederhana):}
\begin{itemize}
  \item Anggota A: kerjakan \texttt{index.html} bagian \texttt{<header>} dan \texttt{<nav>}.
  \item Anggota B: kerjakan \texttt{index.html} bagian \texttt{<main>}.
  \item Anggota C: kerjakan \texttt{index.html} bagian \texttt{<footer>} + buat \texttt{css/style.css}.
\end{itemize}

\begin{tangkapanlayar}
  {assets/images/minggu-15/struktur-branch-github.png}
  {Struktur branch di GitHub --- terlihat branch per anggota.}
  {Buka \emph{repository} di GitHub, klik dropdown branch untuk memastikan semua branch anggota sudah ada.}
\end{tangkapanlayar}

\begin{tangkapanlayar}
  {assets/images/minggu-15/html-semantic-tags.png}
  {Struktur HTML di VS Code dengan semantic tags.}
  {Buka \texttt{index.html} di VS Code, periksa penggunaan tag semantik.}
\end{tangkapanlayar}

\subsection{Styling dengan CSS}
\label{subsec:css-styling}

Buka \texttt{css/style.css} dan mulai \emph{styling} berdasarkan \emph{mockup} Figma.

\begin{lstlisting}[language=CSS, caption={CSS dasar --- reset, header, hero, main, footer, responsive}, label={lst:css-dasar}]
/* Reset & Base Styles */
* {
  margin: 0;
  padding: 0;
  box-sizing: border-box;
}

body {
  font-family: 'Segoe UI', Tahoma, Geneva,
               Verdana, sans-serif;
  line-height: 1.6;
  color: #333;
}

/* Header & Navigation */
header {
  background-color: #2c3e50;
  color: white;
  padding: 1rem 2rem;
}

nav {
  display: flex;
  justify-content: space-between;
  align-items: center;
  max-width: 1200px;
  margin: 0 auto;
}

/* Hero Section */
.hero {
  background-color: #3498db;
  color: white;
  text-align: center;
  padding: 4rem 2rem;
}

/* Main Content */
main {
  max-width: 1200px;
  margin: 0 auto;
  padding: 2rem;
}

/* Footer */
footer {
  background-color: #2c3e50;
  color: white;
  text-align: center;
  padding: 2rem;
  margin-top: 4rem;
}

/* Responsive Design */
@media (max-width: 768px) {
  nav {
    flex-direction: column;
    gap: 1rem;
  }
  .hero {
    padding: 2rem 1rem;
  }
}
\end{lstlisting}

\begin{tip}
Selalu referensikan \emph{mockup} Figma saat menulis CSS --- jangan ``nebak'' ukuran. Gunakan \emph{flexbox} untuk \emph{layout} dan \emph{media query} untuk responsif. Jangan lupa \texttt{box-sizing: border-box} di \emph{reset} --- ini mencegah masalah \emph{layout} umum.
\end{tip}

\subsection{Interaktivitas JavaScript (Opsional)}
\label{subsec:javascript}

Jika \emph{mockup} memerlukan elemen interaktif (misal: menu \emph{hamburger}, \emph{accordion} FAQ, tab navigasi):

\begin{lstlisting}[language=JavaScript, caption={Toggle menu mobile (vanilla JS)}, label={lst:js-toggle}]
// js/script.js

// Contoh: Toggle menu mobile
const menuToggle =
  document.querySelector('.menu-toggle');
const navMenu =
  document.querySelector('.nav-menu');

if (menuToggle) {
  menuToggle.addEventListener('click', () => {
    navMenu.classList.toggle('active');
  });
}
\end{lstlisting}

\begin{tangkapanlayar}
  {assets/images/minggu-15/website-desktop-browser.png}
  {Website tampil di browser (Live Server) --- versi desktop.}
  {Buka \texttt{index.html} melalui Live Server di VS Code, pastikan tampilan sesuai \emph{mockup}.}
\end{tangkapanlayar}

\begin{tangkapanlayar}
  {assets/images/minggu-15/website-mobile-browser.png}
  {Website tampil di browser --- versi mobile (toggle responsive).}
  {Gunakan DevTools $\rightarrow$ \emph{Toggle Device Toolbar} untuk memeriksa tampilan mobile.}
\end{tangkapanlayar}

% ---------------------------------------------------------------------
% 6. SESI CLINIC / MENTORING
% ---------------------------------------------------------------------
\section{Sesi Clinic / Mentoring}
\label{sec:clinic}

Pengajar berkeliling ke setiap kelompok selama sesi ini. Manfaatkan waktu untuk:
\begin{itemize}
  \item review cepat kode dan struktur \emph{repo};
  \item tanya-jawab tentang bagian kode yang belum dipahami;
  \item perbaikan masalah teknis yang ditemukan.
\end{itemize}

\begin{catatan}
Pastikan setiap kelompok bisa menjawab pertanyaan berikut:
\begin{enumerate}
  \item Bagaimana cara website ini mengatasi \emph{pain point} \#1 dari hasil riset?
  \item Siapa yang mengerjakan bagian apa? Apakah semua anggota sudah punya tanggung jawab?
  \item Apakah website ini responsif di mobile? Bisa didemo?
  \item Apakah ada penggunaan \emph{agentic coding tools}? Jika ya, bagian mana?
\end{enumerate}
\end{catatan}

% ---------------------------------------------------------------------
% 7. DEPLOYMENT KE GITHUB PAGES
% ---------------------------------------------------------------------
\section{Deployment ke GitHub Pages}
\label{sec:deployment}

\subsection{Persiapan Repository}
\label{subsec:persiapan-deploy}

\begin{langkah}
  \item Pastikan file sudah ter-\emph{commit}:
\begin{lstlisting}[language=bash, caption={Commit implementasi}]
git add .
git commit -m "feat: implementasi halaman utama sesuai mockup Figma"
git push
\end{lstlisting}
  \item Pastikan \emph{branch} \texttt{main} atau \texttt{final-project} berisi kode terbaru. Jika kode di \emph{branch} lain, \emph{merge} terlebih dahulu:
\begin{lstlisting}[language=bash, caption={Merge ke main}]
git checkout main
git merge final-project
git push origin main
\end{lstlisting}
\end{langkah}

\subsection{Aktifkan GitHub Pages}
\label{subsec:aktifkan-gh-pages}

\begin{langkah}
  \item Buka \emph{repository} di GitHub.
  \item Klik tab \textbf{Settings} $\rightarrow$ geser ke bagian kiri $\rightarrow$ klik \textbf{Pages}.
  \item Di bagian \textbf{Source}:
  \begin{itemize}
    \item Pilih \textbf{``Deploy from a branch''}.
    \item Pilih \emph{branch}: \texttt{main} (atau \texttt{final-project}).
    \item Pilih folder: \texttt{/ (root)} atau \texttt{/docs} (tergantung struktur).
    \item Klik \textbf{Save}.
  \end{itemize}
  \item Tunggu 1--2 menit, lalu \emph{refresh} halaman Settings $\rightarrow$ Pages.
  \item Klik tautan yang muncul di bagian atas (contoh: \texttt{https://username.github.io/repo-name/}).
\end{langkah}

\begin{tangkapanlayar}
  {assets/images/minggu-15/github-pages-settings.png}
  {GitHub Settings $\rightarrow$ Pages --- konfigurasi source.}
  {Buka Settings $\rightarrow$ Pages di \emph{repository} GitHub, pastikan \emph{branch} dan folder sudah benar.}
\end{tangkapanlayar}

\begin{tangkapanlayar}
  {assets/images/minggu-15/website-deployed-browser.png}
  {Website yang sudah ter-deploy di browser --- URL GitHub Pages terlihat.}
  {Buka tautan GitHub Pages di browser, pastikan halaman utama tampil dengan benar.}
\end{tangkapanlayar}

\begin{tangkapanlayar}
  {assets/images/minggu-15/website-deployed-mobile.png}
  {Website di mobile view setelah deploy.}
  {Gunakan DevTools untuk memeriksa tampilan mobile setelah deployment.}
\end{tangkapanlayar}

\subsection{Update Deployment}
\label{subsec:update-deploy}

Setiap kali ada perubahan baru, cukup \emph{push} ke \emph{branch} yang sudah dikonfigurasi:

\begin{lstlisting}[language=bash, caption={Push untuk update otomatis}]
git add .
git commit -m "fix: perbaikan layout header"
git push
\end{lstlisting}

GitHub Pages akan otomatis \emph{update} dalam $\pm$ 1--2 menit.

% ---------------------------------------------------------------------
% 8. TROUBLESHOOTING DEPLOYMENT
% ---------------------------------------------------------------------
\section{Troubleshooting Deployment}
\label{sec:troubleshooting-deploy}

\begin{table}[htbp]
\centering
\caption{Masalah Umum Deployment --- GitHub Pages}
\label{tab:troubleshooting-deploy}
\small
\begin{tabular}{@{}p{5.5cm}p{8cm}@{}}
\toprule
\textbf{Masalah} & \textbf{Solusi} \\
\midrule
Website tidak muncul setelah deploy & Tunggu 2--3 menit. Jika masih tidak muncul, cek tab \textbf{Actions} di \emph{repo} --- lihat apakah \emph{build} berhasil. \\
Halaman 404 Not Found & Pastikan file utama bernama \texttt{index.html} dan berada di \emph{folder root} (atau \emph{folder} yang dipilih di Settings). \\
CSS/JS tidak termuat & Cek path file di HTML --- gunakan \emph{relative path} (\texttt{css/style.css}), bukan \emph{absolute path} (\texttt{/css/style.css}). \\
Gambar tidak muncul & Pastikan file gambar ada di \emph{repo} dan path-nya benar. Perhatikan huruf besar/kecil (\emph{case-sensitive} di GitHub Pages). \\
Deploy dari \emph{branch} selain \texttt{main} & Di Settings $\rightarrow$ Pages, ganti \emph{branch source} ke \emph{branch} yang benar. \\
Subfolder tidak terbaca & Pastikan \texttt{index.html} ada di \emph{root folder} yang dipilih di Settings. \\
\bottomrule
\end{tabular}
\end{table}

% ---------------------------------------------------------------------
% 9. TROUBLESHOOTING UMUM
% ---------------------------------------------------------------------
\section{Troubleshooting Umum}
\label{sec:troubleshooting-umum}

\begin{table}[htbp]
\centering
\caption{Masalah Umum Implementasi dan Solusi}
\label{tab:troubleshooting-umum}
\small
\begin{tabular}{@{}p{5.5cm}p{8cm}@{}}
\toprule
\textbf{Masalah} & \textbf{Solusi} \\
\midrule
CSS tidak ter-apply & Cek \texttt{<link>} di HTML --- pastikan path benar. Cek juga apakah file CSS benar-benar ada di folder yang tepat. \\
Gambar tidak tampil & Gunakan path relatif: \texttt{assets/gambar.png}. Pastikan nama file tepat (\emph{case-sensitive}). \\
Layout pecah di mobile & Pastikan \texttt{<meta name="viewport">} ada di \texttt{<head>}. Cek \emph{media query} di CSS. Gunakan DevTools $\rightarrow$ \emph{Toggle Device Toolbar}. \\
JavaScript tidak jalan & Buka Console di DevTools (F12 $\rightarrow$ Console). Cek \emph{error message}. Pastikan file JS sudah di-\emph{link} di HTML sebelum \texttt{</body>}. \\
Git \emph{merge conflict} & Buka file yang conflict $\rightarrow$ cari \texttt{<<<<<<<} $\rightarrow$ pilih versi yang benar $\rightarrow$ hapus marker conflict $\rightarrow$ \emph{commit}. \\
Flexbox tidak bekerja & Pastikan \emph{parent element} punya \texttt{display: flex}. Cek apakah ada \texttt{width} atau \texttt{height} yang membatasi. \\
\bottomrule
\end{tabular}
\end{table}

% ---------------------------------------------------------------------
% 10. PERSIAPAN PRESENTASI
% ---------------------------------------------------------------------
\section{Persiapan Presentasi}
\label{sec:persiapan-presentasi}

\subsection{Struktur Presentasi}
\label{subsec:struktur-presentasi}

Buat file \texttt{docs/presentasi-outline.md} dengan struktur berikut:
\begin{enumerate}
  \item \textbf{Pembukaan (1 menit)} --- Perkenalan kelompok \& nama website; URL asli vs URL redesign.
  \item \textbf{Riset UI/UX (3 menit)} --- \emph{Screenshot} halaman asli; top 3 \emph{pain point}; temuan \emph{benchmarking}.
  \item \textbf{Desain Ulang (3 menit)} --- \emph{Screenshot mockup} Figma; 3 keputusan desain utama beserta kaitannya dengan temuan riset.
  \item \textbf{Implementasi (3 menit)} --- \emph{Live demo} website; teknologi yang digunakan; demo responsif.
  \item \textbf{Pembagian Kerja \& Penggunaan AI (2 menit)} --- Tabel kontribusi; dokumentasi penggunaan \emph{agentic coding tools}.
  \item \textbf{Penutup (1 menit)} --- Refleksi dan Q\&A.
\end{enumerate}

\subsection{Bagikan Peran Presentasi}
\label{subsec:peran-presentasi}

\begin{itemize}
  \item Anggota 1: Pembukaan \& Riset.
  \item Anggota 2: Desain Ulang.
  \item Anggota 3: Implementasi \& Demo.
  \item Pembagian Kerja \& Penutup: diskusikan sendiri.
\end{itemize}

\begin{tangkapanlayar}
  {assets/images/minggu-15/outline-presentasi.png}
  {Outline presentasi di Markdown.}
  {Buka \texttt{docs/presentasi-outline.md} di VS Code, pastikan seluruh segmen sudah terisi.}
\end{tangkapanlayar}

% ---------------------------------------------------------------------
% 11. FINAL COMMIT \& DOKUMENTATION
% ---------------------------------------------------------------------
\section{Final Commit dan Documentation}
\label{sec:final-commit}

\begin{langkah}
  \item Pastikan \texttt{README.md} di \emph{root repo} sudah lengkap --- nama website, anggota kelompok (beserta NIM dan tugas), \emph{tech stack}, cara menjalankan, tautan deploy, tautan Figma, dan dokumentasi penggunaan \emph{agentic coding tools} (jika ada).
  \item \emph{Commit} dan \emph{push}:
\begin{lstlisting}[language=bash, caption={Final commit documentation}]
git add .
git commit -m "docs: lengkapi README dan siapkan bahan presentasi"
git push
\end{lstlisting}
\end{langkah}

\begin{tangkapanlayar}
  {assets/images/minggu-15/readme-lengkap-github.png}
  {README.md yang sudah lengkap di GitHub.}
  {Buka \emph{repository} di GitHub, pastikan \texttt{README.md} menampilkan seluruh informasi yang diperlukan.}
\end{tangkapanlayar}

\begin{tangkapanlayar}
  {assets/images/minggu-15/commit-history-contributors.png}
  {Commit history --- terlihat kontribusi per anggota.}
  {Buka tab Insights $\rightarrow$ Contributors di \emph{repository} GitHub untuk memastikan kontribusi merata.}
\end{tangkapanlayar}

% ---------------------------------------------------------------------
% 12. PANDUAN AGENTIC CODING TOOLS (OPSIONAL)
% ---------------------------------------------------------------------
\section{Panduan Agentic Coding Tools (Opsional)}
\label{sec:agentic-coding}

Jika kelompok menggunakan GitHub Copilot, Claude Code, atau tools serupa, dokumentasikan penggunaannya.

\begin{catatan}
Tips penggunaan yang bertanggung jawab:
\begin{enumerate}
  \item Selalu \emph{review} kode yang dihasilkan --- jangan langsung \emph{copy-paste}.
  \item Pahami setiap baris --- siapkan diri untuk menjelaskan saat sesi tanya-jawab minggu 16.
  \item Gunakan sebagai bantuan, bukan pengganti --- anggota harus bisa memodifikasi kode sendiri.
  \item Dokumentasikan semua \emph{prompt} di \texttt{README.md} atau file terpisah.
\end{enumerate}
\end{catatan}

% ---------------------------------------------------------------------
% 13. OUTPUT YANG DIHARAPKAN
% ---------------------------------------------------------------------
\section{Output yang Diharapkan}
\label{sec:output-minggu-15}

Setelah sesi ini, setiap kelompok harus memiliki:
\begin{itemize}
  \item Halaman HTML/CSS/JS utama (\texttt{src/index.html}, \texttt{src/css/}, \texttt{src/js/}).
  \item Website ter-deploy (tautan GitHub Pages).
  \item Responsif di mobile \& desktop (sudah diuji di browser).
  \item \texttt{README.md} lengkap di \emph{root repository}.
  \item \emph{Commit history} per anggota (terlihat di GitHub $\rightarrow$ Insights $\rightarrow$ Contributors).
  \item \emph{Outline} presentasi (\texttt{docs/presentasi-outline.md}).
  \item Semua perubahan ter-\emph{push} ke \emph{branch} \texttt{final-project} atau \texttt{main}.
\end{itemize}

% ---------------------------------------------------------------------
% 14. TAUTAN REFERENSI
% ---------------------------------------------------------------------
\section{Tautan Referensi}
\label{sec:referensi-minggu-15}

\begin{itemize}
  \item GitHub Pages Docs: \url{https://docs.github.com/en/pages}
  \item MDN: HTML Semantics: \url{https://developer.mozilla.org/en-US/docs/Guide/HTML/Sections_and_outlines_of_an_HTML5_document}
  \item CSS Flexbox Guide: \url{https://css-tricks.com/snippets/css/a-guide-to-flexbox/}
  \item CSS Media Queries: \url{https://developer.mozilla.org/en-US/docs/Web/CSS/Media_Queries}
  \item Figma to Code: \url{https://www.figma.com/best-practices/from-design-to-code/}
\end{itemize}

% ---------------------------------------------------------------------
% 15. PERSIAPAN UNTUK MINGGU 16
% ---------------------------------------------------------------------
\section{Persiapan untuk Minggu 16 (Presentasi)}
\label{sec:persiapan-minggu-16}

Sebelum pertemuan minggu depan, pastikan:
\begin{enumerate}
  \item Website sudah ter-deploy dan bisa diakses semua orang.
  \item Semua anggota bisa menjelaskan kode bagian masing-masing.
  \item \emph{Outline} presentasi sudah final dan semua anggota tahu siapa bicara kapan.
  \item Siapkan demo langsung --- buka website di browser, siap untuk \emph{resize} dan navigasi.
  \item Latihan presentasi --- coba presentasikan dalam waktu $\pm$ 10 menit.
\end{enumerate}

% ---------------------------------------------------------------------
% 16. LATIHAN
% ---------------------------------------------------------------------
\section{Latihan Mandiri}
\label{sec:latihan-minggu-15}

\begin{latihan}[Eksplorasi Deployment]
\begin{enumerate}[label=\alph*., leftmargin=2em, itemsep=2pt]
  \item Buat \emph{branch} terpisah di \emph{repository} pribadi, tambahkan perubahan kecil pada file HTML, lalu \emph{merge} ke \texttt{main}. Amati apakah GitHub Pages ter-\emph{update} secara otomatis.
  \item Buka DevTools di browser, resize jendela dari 1440px ke 375px, dan identifikasi setidaknya 2 masalah \emph{layout} pada halaman Anda. Perbaiki menggunakan \emph{media query}.
\end{enumerate}
\end{latihan}
